\documentclass[10pt,a4paper]{amsart}
\usepackage{amsmath,amssymb,amsthm,graphicx}
% Retained preamble for the Stacks Project chapter {CH} (shared project preamble).
% The chapter file pinned-original-source.tex only carries \input{{preamble}} on its first line,
% so no preamble text lives inside the chapter itself.  The Stacks Project uses one single shared
% preamble for every chapter: the snapshot stages it three times (sources/stacks-chow/preamble.tex,
% sources/stacks-etale-cohomology/preamble.tex and focused-seeds/AGP/sources/stacks-snapshot-20261011/
% preamble.tex) and those copies are identical in size.  The block below is copied mechanically,
% line for line, from sources/stacks-chow/preamble.tex lines 14-18 and 154-217: the Xy-pic setup and
% the native theorem declarations together with the mathematical macros.  Nothing was added, removed
% or redefined; the wrapper class and the amsmath/amssymb/amsthm/graphicx base packages are supplied
% by the assembler.
% For commutative diagrams we use Xy-pic
% Chapter: stacks-cohomology.tex -- Cohomology of Sheaves.
\usepackage[all]{xy}

% We use 2cell for 2-commutative diagrams.
\xyoption{2cell}

% Theorem environments.
%
\theoremstyle{plain}
\newtheorem{theorem}[subsection]{Theorem}
\newtheorem{proposition}[subsection]{Proposition}
\newtheorem{lemma}[subsection]{Lemma}

\theoremstyle{definition}
\newtheorem{definition}[subsection]{Definition}
\newtheorem{example}[subsection]{Example}
\newtheorem{exercise}[subsection]{Exercise}
\newtheorem{situation}[subsection]{Situation}

\theoremstyle{remark}
\newtheorem{remark}[subsection]{Remark}
\newtheorem{remarks}[subsection]{Remarks}

\numberwithin{equation}{subsection}

% Macros
%
\def\lim{\mathop{\mathrm{lim}}\nolimits}
\def\colim{\mathop{\mathrm{colim}}\nolimits}
\def\Spec{\mathop{\mathrm{Spec}}}
\def\Hom{\mathop{\mathrm{Hom}}\nolimits}
\def\Ext{\mathop{\mathrm{Ext}}\nolimits}
\def\SheafHom{\mathop{\mathcal{H}\!\mathit{om}}\nolimits}
\def\SheafExt{\mathop{\mathcal{E}\!\mathit{xt}}\nolimits}
\def\Sch{\mathit{Sch}}
\def\Mor{\mathop{\mathrm{Mor}}\nolimits}
\def\Ob{\mathop{\mathrm{Ob}}\nolimits}
\def\Sh{\mathop{\mathit{Sh}}\nolimits}
\def\NL{\mathop{N\!L}\nolimits}
\def\CH{\mathop{\mathrm{CH}}\nolimits}
\def\proetale{{pro\text{-}\acute{e}tale}}
\def\etale{{\acute{e}tale}}
\def\QCoh{\mathit{QCoh}}
\def\Ker{\mathop{\mathrm{Ker}}}
\def\Im{\mathop{\mathrm{Im}}}
\def\Coker{\mathop{\mathrm{Coker}}}
\def\Coim{\mathop{\mathrm{Coim}}}

% Boxtimes
%
\DeclareMathSymbol{\boxtimes}{\mathbin}{AMSa}{"02}

%
% Macros for moduli stacks/spaces
%
\def\QCohstack{\mathcal{QC}\!\mathit{oh}}
\def\Cohstack{\mathcal{C}\!\mathit{oh}}
\def\Spacesstack{\mathcal{S}\!\mathit{paces}}
\def\Quotfunctor{\mathrm{Quot}}
\def\Hilbfunctor{\mathrm{Hilb}}
\def\Curvesstack{\mathcal{C}\!\mathit{urves}}
\def\Polarizedstack{\mathcal{P}\!\mathit{olarized}}
\def\Complexesstack{\mathcal{C}\!\mathit{omplexes}}
% \Pic is the operator that assigns to X its picard group, usage \Pic(X)
% \Picardstack_{X/B} denotes the Picard stack of X over B
% \Picardfunctor_{X/B} denotes the Picard functor of X over B
\def\Pic{\mathop{\mathrm{Pic}}\nolimits}
\def\Picardstack{\mathcal{P}\!\mathit{ic}}
\def\Picardfunctor{\mathrm{Pic}}
\def\Deformationcategory{\mathcal{D}\!\mathit{ef}}
\usepackage{cleveref}
\crefname{originallemma}{original lemma}{original lemmas}
\Crefname{originallemma}{Original lemma}{Original lemmas}
\makeatletter
\newlabel{lemma-mayer-vietoris}{{01EB}{source}{}{sourcepointer-1}{}}
\newlabel{lemma-mayer-vietoris@cref}{{[originallemma][0][]01EB}{[1][0][]source}{}{}{}}
\makeatother
\crefname{originallemma}{original lemma}{original lemmas}
\Crefname{originallemma}{Original lemma}{Original lemmas}
\makeatletter
\newlabel{lemma-cohomology-of-open}{{01E1}{source}{}{sourcepointer-2}{}}
\newlabel{lemma-cohomology-of-open@cref}{{[originallemma][0][]01E1}{[1][0][]source}{}{}{}}
\makeatother
\title{Relative Mayer-Vietoris long exact sequence for higher direct images of a two-open cover}
\author{The Stacks Project Authors}
\date{stacks-project:cohomology.tex:local-snapshot-20261011; GNU FDL 1.2 or later}
\begin{document}
\maketitle
\noindent\textit{Editorial note.} All delivered mathematical text is the author's own, copied from the pinned Stacks Project chapter file cohomology.tex; no statement or proof step was rewritten, reordered or removed. The author carries the functoriality claim over from the absolute statement without proof, and the bracketed name Relative Mayer-Vietoris in the environment header is reproduced as written.\par\medskip
\noindent\textit{Editorial note.} The chapter file pinned-original-source.tex carries only \textbackslash{}input\{preamble\} on its first line, so no preamble text lives inside the chapter. The retained file is a mechanical, line-for-line copy of the project's single shared preamble as staged in the snapshot: sources/stacks-chow/preamble.tex lines 14-18 (Xy-pic setup) and 154-217 (theorem declarations and mathematical macros). Nothing was added, removed or redefined, and the excerpt serves only as preamble, never as a span coordinate.\par\medskip
\noindent\textit{Editorial note.} Both referenced results live in the same chapter and are located at their \textbackslash{}label lines inside the pinned chapter file.\par\medskip
\noindent\textit{Allowed standard prerequisites (editorial).}\begin{itemize}
\item Mayer-Vietoris for a two-open cover, in its degreewise form. Same chapter, Lemma labelled lemma-mayer-vietoris (source line 767, Stacks tag 01EB); invoked at proof line 842, and its proof is where the displayed degreewise short exact sequence was established. Supplies the degreewise exactness of the sequence of sections over f\textasciicircum{}\{-1\}(W), which is what makes the sequence of complexes of direct images exact.
\item Cohomology of an open subspace is computed by restriction of an injective resolution. Same chapter, Lemma labelled lemma-cohomology-of-open (source line 496, Stacks tag 01E1); invoked at proof lines 844-846. Supplies that the restrictions of an injective resolution are injective resolutions of the restricted modules, which identifies the cohomology sheaves appearing in the long exact sequence with the higher direct images.
\end{itemize}\medskip
\ifdefined\hypertarget\hypertarget{sourcepointer-1}{}\fi
\noindent\textit{Source reference.} Original lemma 01EB: sources/stacks-cohomology/pinned-original-source.tex, source lines 767--767. Same-chapter pointer invoked at proof line 842 for the degreewise short exact sequence of sections; the assembly of the sequence of complexes of direct images and the reading off of the long exact sequence are retained in the counted span.\par\medskip
\ifdefined\hypertarget\hypertarget{sourcepointer-2}{}\fi
\noindent\textit{Source reference.} Original lemma 01E1: sources/stacks-cohomology/pinned-original-source.tex, source lines 496--496. Same-chapter pointer invoked at proof lines 844-846 to identify the restrictions of the injective resolution as injective resolutions; the passage from the exact sequence of complexes to the long exact sequence of higher direct images is retained in the counted span.\par\medskip
\begin{lemma}[Relative Mayer-Vietoris]
\label{lemma-relative-mayer-vietoris}
Let $f : X \to Y$ be a morphism of ringed spaces.
Suppose that $X = U \cup V$ is a union of two open subsets.
Denote $a = f|_U : U \to Y$, $b = f|_V : V \to Y$, and
$c = f|_{U \cap V} : U \cap V \to Y$.
For every $\mathcal{O}_X$-module $\mathcal{F}$
there exists a long exact sequence
$$
0 \to
f_*\mathcal{F} \to
a_*(\mathcal{F}|_U) \oplus b_*(\mathcal{F}|_V) \to
c_*(\mathcal{F}|_{U \cap V}) \to
R^1f_*\mathcal{F} \to \ldots
$$
This long exact sequence is functorial in $\mathcal{F}$.
\end{lemma}

\begin{proof}
Let $\mathcal{F} \to \mathcal{I}^\bullet$ be an injective resolution
of $\mathcal{F}$. We claim that we
get a short exact sequence of complexes
$$
0 \to
f_*\mathcal{I}^\bullet \to
a_*\mathcal{I}^\bullet|_U \oplus b_*\mathcal{I}^\bullet|_V \to
c_*\mathcal{I}^\bullet|_{U \cap V} \to
0.
$$
Namely, for any open $W \subset Y$, and for any $n \geq 0$ the
corresponding sequence of groups of sections over $W$
$$
0 \to
\mathcal{I}^n(f^{-1}(W)) \to
\mathcal{I}^n(U \cap f^{-1}(W))
\oplus \mathcal{I}^n(V \cap f^{-1}(W)) \to
\mathcal{I}^n(U \cap V \cap f^{-1}(W)) \to
0
$$
was shown to be short exact in the proof of Lemma \ref{lemma-mayer-vietoris}.
The lemma follows by taking cohomology sheaves and using the fact that
$\mathcal{I}^\bullet|_U$ is an injective resolution of $\mathcal{F}|_U$
and similarly for $\mathcal{I}^\bullet|_V$, $\mathcal{I}^\bullet|_{U \cap V}$
see Lemma \ref{lemma-cohomology-of-open}.
\end{proof}

\end{document}
